\title{Group Sequence Policy Optimization}

\begin{abstract}
We present Group Sequence Policy Optimization (GSPO), a robust, effective, and high-performing reinforcement learning method tailored for training large language models. Distinct from earlier approaches that rely on token-level importance ratios, GSPO employs a sequence-level perspective by basing the importance ratio on entire sequence likelihoods, alongside applying clipping, reward assignment, and optimization at the sequence level. Our results show that GSPO outperforms the GRPO algorithm in terms of both training efficiency and overall effectiveness, significantly enhances the stability of Mixture-of-Experts (MoE) reinforcement learning, and can streamline RL infrastructure development. The strengths of GSPO have substantially contributed to the notable progress seen in the recent Qwen3 models.
\end{abstract}

\section{Introduction}
\label{sec:intro}

Reinforcement learning (RL) has become a central framework in advancing the scalability of language models \citep{o1,dpsk-r1,qwq32b,qwen3}. When applied at scale, RL enables language models to address complex challenges—including tasks found in competitive mathematics and computer programming—by engaging in more elaborate and extended reasoning sequences.

Achieving effective scaling of RL through increased computational resources fundamentally depends on preserving consistent and resilient training behavior. Yet, leading RL approaches such as GRPO \citep{grpo} face significant challenges in this regard, displaying pronounced instability when applied to the training of very large language models. This often culminates in complete and irrecoverable model failures \citep{qwen3, minimax-m1}. Such instability significantly constrains attempts to extend the capacities of language models via sustained RL optimization.

In this work, we demonstrate that the root cause of GRPO’s instability lies in the incorrect use and breakdown of importance sampling weights within its algorithm. This flaw results in significant high-variance noise during training, which grows as response length increases and is exacerbated by the clipping operation. Collectively, these issues drive the model towards collapse.

To overcome these fundamental challenges, we introduce \textbf{Group Sequence Policy Optimization (GSPO)}, a novel reinforcement learning approach tailored for the training of large language models. The principal contribution of GSPO is its rigorous formulation of the importance ratio, which is derived from sequence likelihood \citep{click} and is consistent with the foundational idea of importance sampling. Furthermore, GSPO evaluates normalized rewards as the advantages across several responses for each query, thereby promoting consistency between sequence-based reward assignment and the optimization process.

Our experimental results highlight GSPO’s marked advantages over GRPO in terms of training stability, efficiency, and overall effectiveness. Importantly, GSPO naturally addresses the instability issues associated with reinforcement learning for large Mixture-of-Experts (MoE) models, obviating the reliance on intricate stabilization mechanisms and suggesting opportunities to streamline RL infrastructure. These strengths of GSPO have led directly to substantial performance gains in the newest Qwen3 models. We anticipate that GSPO will serve as a dependable and extensible algorithmic base to drive ongoing progress in large-scale RL training for language models.

\section{Preliminaries}

\paragraph{Notation}
Throughout this work, we represent an autoregressive language model, parameterized by $\theta$, as a policy $\pi_\theta$. The variable $x$ stands for a query, and the collection of all queries is denoted by $\mathcal{D}$. For a query $x$ with corresponding response $y$, the probability of $y$ under the policy $\pi_\theta$ is written as $\pi_\theta (y | x)=\prod_{t=1}^{|y|} \pi_\theta (y_t | x, y_{<t} )$, where $|y|$ refers to the total number of tokens in the sequence $y$. Evaluation of a query-response pair $(x, y)$ is carried out using a verifier $r$, assigning a reward $r(x, y) \in [0, 1]$.

\paragraph{Proximal Policy Optimization (PPO)} Proximal Policy Optimization (PPO) \citep{ppo} updates policies using data collected from the previous policy $\pi_{\theta_\text{old}}$, restricting changes to remain close to this reference through a clipping method. More precisely, PPO maximizes the following surrogate objective during policy training (the KL regularization component is omitted hereafter for conciseness, since it is not central to our analysis):

\begin{align}
\mathcal{J}_\text{PPO}(\theta) = \mathbb{E}_{ x \sim \mathcal{D},\, y \sim \pi_{\theta_\text{old}}( \cdot | x) }
\left[ \frac{1}{|y|} \sum_{t=1}^{|y|} 
\min \left( w_{t}(\theta) \widehat{A}_{t},  \, \mathrm{clip} \left( w_{t}(\theta), 1 - {\varepsilon}, 1 + {\varepsilon}\right) \widehat{A}_{t} \right)
\right],
\end{align}

Here, the token importance ratio $w_{t}(\theta)$ is specified as $
w_{t}(\theta) = \frac{ \pi_{\theta} (y_{t} | x, y_{<t}) }{ \pi_{\theta_\text{old}} (y_{t} | x,y_{<t})} $. The advantage term $\widehat{A}_{t}$ associated with $y_t$ is computed using a separate value model, and the parameter $\varepsilon$ denotes the allowed clipping interval for the importance ratios.

A fundamental practical limitation of PPO stems from its significant dependence on the value model. Notably, the value model typically matches the policy model in scale, resulting in substantial demands on memory and computation. In addition, the success of the algorithm is contingent upon having an accurate value estimate. Developing a robust value model is intrinsically difficult, and achieving scalability for longer outputs and more intricate tasks proves to be even more demanding.

\paragraph{Group Relative Policy Optimization (GRPO)} GRPO \citep{grpo} eliminates the dependence on a value model by evaluating the comparative advantage of individual responses within a set of responses to an identical query. In particular, GRPO seeks to maximize the following objective:

\begin{align}
\label{equ:grpo}
\mathcal{J}_\text{GRPO}(\theta) = \mathbb{E}_{ x \sim \mathcal{D},\, \{y_i\}_{i=1}^G \sim \pi_{\theta_\text{old}}( \cdot | x) }
\left[ \frac{1}{G} \sum_{i=1}^{G} \frac{1}{|y_i|} \sum_{t=1}^{|y_i|} 
\min \left( w_{i,t}(\theta) \widehat{A}_{i,t},  \, \mathrm{clip} \left( w_{i,t}(\theta), 1 - {\varepsilon}, 1 + {\varepsilon}\right) \widehat{A}_{i,t} \right)
\right],
\end{align}

Here, $G$ denotes the total number of responses produced for every query $x$ (also referred to as the group size), and the expressions $w_{i,t}(\theta)$ and $\widehat{A}_{i,t}$ represent, respectively, the importance ratio and the advantage for the token $y_{i,t}$ as follows:

\begin{align}
    w_{i,t}(\theta)=\frac{ \pi_{\theta} (y_{i,t} | x, y_{i,<t}) }{ \pi_{\theta_\text{old}} (y_{i,t} | x,y_{i,<t})},\quad\ 
    \widehat{A}_{i,t} = \widehat{A}_{i} = \frac{r(x, y_i) - \mathrm{mean} \left( \{ r(x, y_i) \}_{i=1}^G \right) }{ \mathrm{std} \left( \{ r(x, y_i) \}_{i=1}^G \right) },
\end{align}

respectively; each token in $y_i$ is therefore assigned the identical advantage value $\widehat{A}_{i}$.

\section{Motivation}
\label{sec:motivation}

As models increase in size, utilize sparse architectures such as Mixture-of-Experts, and generate longer responses, leveraging large rollout batch sizes becomes essential for achieving optimal hardware throughput in RL. To enhance sample efficiency, a common approach involves dividing the overall batch of rollout experiences into several mini-batches for gradient computation. However, this strategy introduces an inherent off-policy aspect, since the responses $y$ are drawn from a previous policy $\pi_{\theta_\text{old}}$ instead of the current policy $\pi_\theta$ being updated. This situation underscores why clipping mechanisms, like those used in PPO and GRPO, are employed—to restrict the influence of data points that deviate substantially from the present policy and thus mitigate the effects of off-policy learning on gradient calculations.

Although techniques such as clipping seek to address off-policy discrepancies, we observe a deeper problem inherent to GRPO: \textit{the objective itself is not well-defined}. This issue is exacerbated during the training of large models on tasks that require extended responses, often resulting in severe model degradation. The root cause lies in the improper use of importance sampling weights within the GRPO objective. In essence, importance sampling is meant to approximate the expectation of a function $f$ under a target distribution $\pi_\text{tar}$ by applying suitable weights to samples collected from a behavior distribution $\pi_\text{beh}$:

\begin{align}
\mathbb{E}_{z \sim \pi_\text{tar}} \left[ f(z) \right] = \mathbb{E}_{z \sim \pi_\text{beh}} \left[ \frac{ \pi_\text{tar} (z) }{ \pi_\text{beh} (z)} \, f(z) \right].
\end{align}

Importantly, this approach depends on taking the average across a large number of samples ($N \gg 1$) drawn from the behavior distribution $\pi_\text{beh}$ so that the importance weight $\frac{ \pi_\text{tar} (z) }{ \pi_\text{beh} (z)}$ can adequately adjust for discrepancies between the distributions.

By contrast, GRPO incorporates the importance weight $\frac{ \pi_{\theta} (y_{i,t} | x, y_{i,<t}) }{ \pi_{\theta_\text{old}} (y_{i,t} | x,y_{i,<t})}$ at every token timestep $t$. Because this factor is computed from just a single token sample $y_{i,t}$ taken from the respective next-token distribution $\pi_{\theta_\text{old}}( \cdot | x, y_{i, <t})$, it does not adequately correct for discrepancies between distributions. Rather, it contributes substantial variance to the gradient estimates, with this noisy signal amplifying as sequence length increases and becoming worse when clipping is used. Our experiments show that this sometimes precipitates catastrophic model collapse that typically cannot be undone. When such collapse happens, further training remains ineffective—restarting from an earlier checkpoint, extensively adjusting hyperparameters (such as clipping thresholds), increasing the sequence length, or modifying RL queries all fail to restore the model.

This observation highlights a key limitation inherent in GRPO’s framework. Specifically, the ineffectiveness of importance weighting at the token level reveals an essential insight: \textbf{the granularity of the optimization target must align with the granularity of the reward signal}. Given that rewards are assigned to the full sequence rather than individual tokens, implementing off-policy correction for tokens introduces difficulties. Consequently, we are led to move away from objectives defined at the token level, and instead, investigate approaches that incorporate importance weighting and conduct optimization over the \textit{sequence level} directly.

\section{Algorithm}

\subsection{GSPO: Group Sequence Policy Optimization}

Although the token-level importance weight $\frac{ \pi_{\theta} (y_{i,t} | x, y_{i,<t}) }{ \pi_{\theta_\text{old}} (y_{i,t} | x,y_{i,<t})}$ presents challenges within GRPO, we note that, in language generation scenarios, the \textit{sequence-level} importance weight $\frac{ \pi_{\theta} (y | x) }{ \pi_{\theta_\text{old}} (y | x)}$ holds clear theoretical significance. Specifically, it quantifies the extent to which a response $y$—generated by sampling from $\pi_{\theta_\text{old}} (\cdot | x)$—differs from the distribution $\pi_{\theta} (\cdot | x)$. This alignment is congruent with the sequence-level reward structure and offers an intuitive and meaningful measure for the clipping process.

Drawing upon this simple insight, we introduce the \textbf{Group Sequence Policy Optimization (GSPO)} algorithm. GSPO utilizes the subsequent optimization objective at the sequence level:

\begin{align}
\mathcal{J}_\text{GSPO} (\theta) =
\mathbb{E}_{ x \sim \mathcal{D},\, \{y_i\}_{i=1}^G \sim \pi_{\theta_\text{old}}( \cdot | x) }
\left[ 
\frac{1}{G} \sum_{i=1}^{G}
\min \left( s_{i}(\theta)  \widehat{A}_{i},  \, \mathrm{clip} \left( s_{i}(\theta), 1 - {\varepsilon}, 1 + {\varepsilon} \right) \widehat{A}_{i} \right) 
\right],
\label{equ:gspo}
\end{align}

in which we utilize group-based estimation of advantages:

\begin{align}
\widehat{A}_{i} = \frac{r(x, y_i) - \mathrm{mean} \left( \{ r(x, y_i) \}_{i=1}^G \right) }{ \mathrm{std} \left( \{ r(x, y_i) \}_{i=1}^G \right) },
\end{align}

and establish the importance ratio $s_{i}(\theta)$ utilizing sequence likelihood as described by \citet{click}:

\begin{align}
s_{i}(\theta) = \left( \frac{ \pi_{\theta} (y_i | x) }{ \pi_{\theta_\text{old}} (y_i | x)} \right)^{\frac{1}{|y_i|}}
=
\exp \left( \frac{1}{|y_i|} \sum_{t=1}^{|y_i|} \log \frac{ \pi_{\theta} (y_{i,t} | x, y_{i,<t}) }{ \pi_{\theta_\text{old}} (y_{i,t} | x,y_{i,<t})} \right).
\end{align}

Consequently, GSPO performs clipping at the response level rather than at the token level, thereby removing highly ``off-policy'' samples from the gradient calculation. This approach aligns with both sequence-level reward assignment and the corresponding optimization process. It is important to highlight that length normalization is utilized in $s_{i}(\theta)$ to mitigate variance and keep $s_{i}(\theta)$ within a consistent scale. Without this normalization, even small variations in token likelihoods could cause substantial instability in the sequence-level importance ratio, necessitating distinct clipping thresholds for responses of different lengths. Additionally, the magnitudes of the clipping ranges in GSPO and prior methods such as GRPO often diverge considerably, which stems from the different underlying definitions of the importance ratios in each algorithm.

\subsection{Gradient Analysis}
\label{subsec:gradient}

The gradient of the GSPO objective can be determined in the following manner (for simplicity, the clipping step is not shown):

\begin{align}
\nabla_{\theta} \mathcal{J}_\text{GSPO} (\theta)
=&\ 
\nabla_{\theta} \mathbb{E}_{ x \sim \mathcal{D},\, \{y_i\}_{i=1}^G \sim \pi_{\theta_\text{old}}( \cdot | x) }
\left[ 
\frac{1}{G} \sum_{i=1}^{G} s_{i}(\theta) \widehat{A}_{i}
\right] \\
= &\ 
\mathbb{E}_{ x \sim \mathcal{D},\, \{y_i\}_{i=1}^G \sim \pi_{\theta_\text{old}}( \cdot | x) }
\left[ 
\frac{1}{G} \sum_{i=1}^{G}
s_{i}(\theta) \widehat{A}_{i}
\cdot \nabla_{\theta} \log s_{i}(\theta)
\right] \\
=&\ 
\mathbb{E}_{ x \sim \mathcal{D},\, \{y_i\}_{i=1}^G \sim \pi_{\theta_\text{old}}( \cdot | x) }
\left[ 
\frac{1}{G} \sum_{i=1}^{G} 
\left( \frac{ \pi_{\theta} (y_i | x) }{ \pi_{\theta_\text{old}} (y_i | x)} \right)^{\frac{1}{|y_i|}} \widehat{A}_{i} 
\cdot \frac{1}{|y_i|} \sum_{t=1}^{|y_i|} \nabla_{\theta} \log \pi_{\theta} (y_{i,t} | x, y_{i,<t}) 
\right].
\label{equ:gspo_grad}
\end{align}

For reference, the gradient of the GRPO objective can be expressed as follows (where $\widehat{A}_{i,t}$ is equivalent to $\widehat{A}_{i}$):

\begin{align}
\nabla_{\theta} \mathcal{J}_\text{GRPO} (\theta) 
=&\ 
\nabla_{\theta} \mathbb{E}_{ x \sim \mathcal{D},\, \{y_i\}_{i=1}^G \sim \pi_{\theta_\text{old}}( \cdot | x) }
\left[ \frac{1}{G} \sum_{i=1}^{G} \frac{1}{|y_i|} \sum_{t=1}^{|y_i|} 
w_{i,t}(\theta) \widehat{A}_{i,t}
\right] \\
=&\
\mathbb{E}_{ x \sim \mathcal{D},\, \{y_i\}_{i=1}^G \sim \pi_{\theta_\text{old}}( \cdot | x) }
\left[ \frac{1}{G} \sum_{i=1}^{G} \widehat{A}_{i} 
\cdot \frac{1}{|y_i|} \sum_{t=1}^{|y_i|} 
\frac{ \pi_{\theta} (y_{i,t} | x, y_{i,<t}) }{ \pi_{\theta_\text{old}} (y_{i,t} | x,y_{i,<t})} 
\nabla_{\theta} \log \pi_{\theta} (y_{i,t} | x, y_{i,<t})  
\right].
\end{align}

The key difference between GSPO and GRPO centers on \textit{the method by which they assign weights to the gradients of the tokens’ log likelihoods}. In the case of GRPO, each token is assigned a weight corresponding to its ``importance weight'' $\frac{ \pi_{\theta} (y_{i,t} | x, y_{i,<t}) }{ \pi_{\theta_\text{old}} (y_{i,t} | x,y_{i,<t})}$. These variable weights, ranging over $(0, 1+\varepsilon]$ when $\widehat{A}_{i} >0$ and $[1-\varepsilon, +\infty)$ when $\widehat{A}_{i} < 0$, are significant and not dismissible; as training continues, the accumulation of such disparities can produce erratic effects. Conversely, GSPO assigns uniform weighting across all tokens in a generated response, thereby removing the source of instability inherent to GRPO.

\subsection{GSPO-token: A Token-level Objective Variant}

In contexts such as multi-turn RL, a more detailed advantage modification than what is provided at the sequence level may be necessary. To address this, we present a token-level variant of the GSPO objective, referred to as \textbf{GSPO-token}, which enables the tailoring of advantages on a per-token basis:

\begin{align}
\mathcal{J}_\text{GSPO-token}(\theta)
=
\mathbb{E}_{ x \sim \mathcal{D},\, \{y_i\}_{i=1}^G \sim \pi_{\theta_\text{old}}( \cdot | x) }
\left[ 
\frac{1}{G} \sum_{i=1}^{G} \frac{1}{|y_i|} \sum_{t=1}^{|y_i|} 
\min \left( s_{i,t}(\theta) \widehat{A}_{i,t},  \, \mathrm{clip} \left( s_{i,t}(\theta), 1 - {\varepsilon}, 1 + {\varepsilon} \right) \widehat{A}_{i,t} \right) 
\right],
\label{equ:gspo-token}
\end{align}

in which

\begin{align}
s_{i,t}(\theta) 
= \mathrm{sg} \left[ s_{i}(\theta) \right]  \cdot \frac{ \pi_{\theta} (y_{i,t} | x, y_{i,<t}) }{ \mathrm{sg} \left[ \pi_{\theta} (y_{i,t} | x, y_{i,<t}) \right] },
\end{align}

and $\mathrm{sg}[\cdot]$ indicates extracting only the numerical value while preventing gradient flow, which aligns with the \texttt{detach} function in PyTorch. The gradient for GSPO-token can thus be expressed as:

\begin{align}
\nabla_{\theta} \mathcal{J}_\text{GSPO-token}(\theta)
=&\ 
\nabla_{\theta} \mathbb{E}_{ x \sim \mathcal{D},\, \{y_i\}_{i=1}^G \sim \pi_{\theta_\text{old}}( \cdot | x) }
\left[ 
\frac{1}{G} \sum_{i=1}^{G}
\frac{1}{|y_i|} \sum_{t=1}^{|y_i|} 
s_{i,t}(\theta) \widehat{A}_{i,t}
\right] \\
=&\ 
\mathbb{E}_{ x \sim \mathcal{D},\, \{y_i\}_{i=1}^G \sim \pi_{\theta_\text{old}}( \cdot | x) }
\left[ 
\frac{1}{G} \sum_{i=1}^{G} s_i (\theta) \cdot \frac{1}{|y_i|} \sum_{t=1}^{|y_i|} 
\widehat{A}_{i,t} \frac{ \nabla_{\theta} \pi_{\theta} (y_{i,t} | x, y_{i,<t}) }{ \pi_{\theta} (y_{i,t} | x, y_{i,<t}) }
\right] \\
=&\ 
\mathbb{E}_{ x \sim \mathcal{D},\, \{y_i\}_{i=1}^G \sim \pi_{\theta_\text{old}}( \cdot | x) }
\left[ 
\frac{1}{G} \sum_{i=1}^{G}  \left( \frac{ \pi_{\theta} (y_i | x) }{ \pi_{\theta_\text{old}} (y_i | x)} \right)^{\frac{1}{|y_i|}}
\cdot \frac{1}{|y_i|} \sum_{t=1}^{|y_i|}
\widehat{A}_{i,t} \nabla_{\theta} \log \pi_{\theta} (y_{i,t} | x, y_{i,<t})  
\right].
\label{equ:gspo-token_grad}
\end{align}

Observe that $\frac{ \pi_{\theta} (y_{i,t} | x, y_{i,<t}) }{ \mathrm{sg} \left[ \pi_{\theta} (y_{i,t} | x, y_{i,<t}) \right] }$ evaluates to 1 in practice, which implies $s_{i,t}(\theta)$ and $s_{i}(\theta)$ are numerically equivalent. By comparing Equation~\eqref{equ:gspo} with~\eqref{equ:gspo-token}, as well as Equation~\eqref{equ:gspo_grad} with~\eqref{equ:gspo-token_grad}, it becomes clear that, when the advantage for each token in a given response $y_i$ is set uniformly (specifically, $\widehat{A}_{i,t} = \widehat{A}_{i}$), both GSPO-token and GSPO yield the same numerical outcomes for the optimization target, the clipping criterion, and the theoretical gradient. Nevertheless, GSPO-token provides the additional benefit of allowing the advantage values to differ at the token level, thereby offering greater adaptability.

\section{Experiments and Discussion}

\subsection{Empirical Results}

We conduct experiments using a cold-start model initialized from Qwen3-30B-A3B-Base, and we present both the training reward trajectories and model performance metrics on the AIME'24 (mean Pass@1 across 32 samples), LiveCodeBench (from 202410 to 202502, mean Pass@1 across 8 samples), and CodeForces (Elo Rating) evaluation sets. Throughout the RL training process, each rollout batch is divided into four mini-batches to perform gradient updates. For GSPO, we specify the left and right clipping thresholds in Equation~\eqref{equ:gspo} as 3e-4 and 4e-4, respectively. Our comparison utilizes GRPO as the reference method, with left and right clipping thresholds in Equation~\eqref{equ:grpo} adjusted to 0.2 and 0.27, respectively, following careful parameter tuning for an equitable comparison. It should be emphasized that GRPO requires the Routing Replay training method to achieve stable MoE RL convergence—a point elaborated in \S~\ref{subsec:routing_replay}—whereas \textit{GSPO eliminates the necessity for this approach}.

As illustrated in Figure~\ref{fig:results}, the training process using GSPO maintains consistent stability. Our findings indicate that \textbf{GSPO facilitates ongoing gains in performance by scaling training compute, routinely refreshing the query set, and allowing for longer generation outputs}. Additionally, when compared with GRPO, GSPO exhibits greater training efficiency, attaining higher training accuracy and superior benchmark results given identical amounts of training compute and query consumption. Finally, GSPO has been effectively integrated into RL training for the most recent Qwen3 models, offering strong evidence that GSPO enables RL scaling to significantly enhance large language model capabilities.

\subsection{Curious Observation on Clipping Fractions}

One of the principal differences between GSPO and GRPO lies in the fact that GSPO applies clipping to entire sequences instead of operating at the token level. As illustrated in Figure~\ref{fig:clipping}, GSPO results in a proportion of clipped tokens that is approximately two orders of magnitude higher than GRPO, and this large discrepancy persists regardless of how the clipping range is set. Interestingly, even though GSPO excludes a much higher percentage of tokens—thereby reducing the number used for training or for estimating gradients—it nonetheless demonstrates greater training efficiency than GRPO. This surprising outcome, where more extensive token clipping is associated with improved training efficiency, suggests that GRPO’s token-level gradient estimates are intrinsically noisy and inefficient for leveraging available samples. By contrast, the sequence-level method used in GSPO offers a more stable and effective learning signal.

\subsection{Benefit of GSPO for MoE Training}
\label{subsec:routing_replay}

\paragraph{Background}
MoE models, due to their sparsely activated structure, present distinct stability issues during RL training that are less pronounced in densely activated architectures. Specifically, our investigations reveal that under the GRPO algorithm, MoE models often experience heightened \textit{expert-activation volatility}, which can impede the successful convergence of RL optimization. This is evident in the significant shifts in which experts are selected for the same output response after one or more gradient updates. As an illustration, for the 48-layer Qwen3-30B-A3B-Base model, analysis after each RL gradient update indicates that, for any given rollout sample, approximately 10\% of the experts newly activated under the updated policy $\pi_{\theta}$ do not overlap with those under the previous policy $\pi_{\theta_\text{old}}$. Such instability is exacerbated in deeper MoE architectures, resulting in severe oscillations of the token-level importance weights $w_{i,t}(\theta) = \frac{ \pi_{\theta} (y_{i,t} | x, y_{i,<t}) }{ \pi_{\theta_\text{old}} (y_{i,t} | x,y_{i,<t})}$, thereby rendering these estimates unreliable (as elaborated in \S~\ref{sec:motivation} and~\ref{subsec:gradient}) and ultimately disrupting RL training convergence.

\paragraph{Our Previous Approach} In our earlier work, we addressed this issue by utilizing the \textbf{Routing Replay} training scheme. This approach involves storing the expert selections made by $\pi_{\theta_\text{old}}$ and subsequently reusing these routing decisions within $\pi_{\theta}$ during the computation of importance weights $w_{i,t}(\theta) = \frac{ \pi_{\theta} (y_{i,t} | x, y_{i,<t}) }{ \pi_{\theta_\text{old}} (y_{i,t} | x,y_{i,<t})}$. This procedure ensures that both $\pi_{\theta} (y_{i,t} | x, y_{i,<t})$ and $\pi_{\theta_\text{old}} (y_{i,t} | x,y_{i,<t})$ for each token $y_{i,t}$ activate the identical subset of the network, thereby promoting stability in token-wise importance ratio estimates and facilitating optimization over a consistent expert configuration as parameters are updated. As illustrated in Figure~\ref{fig:routing_replay}, Routing Replay is indispensable for the stable and effective convergence of MoE model training with the GRPO method.

\paragraph{Benefit of GSPO} While Routing Replay ensures that GRPO training for MoE models converges successfully, it introduces extra costs in both memory usage and communication by repeatedly reapplying the same routing patterns, which can ultimately restrict the effective capacity of the MoE model. GSPO, in comparison—and as depicted in Figure~\ref{fig:results}—removes the necessity for Routing Replay altogether. It facilitates the standard computation of importance ratios $s_i(\theta)$, achieving stable optimization and normal convergence. The principal idea behind GSPO is its exclusive reliance on sequence likelihoods (i.e., $\pi_{\theta} (y_i | x)$), rather than focusing on the likelihoods of individual tokens (i.e., $\pi_{\theta} (y_{i,t} | x, y_{i,<t})$). Because the MoE model continually preserves robust language modeling abilities, the sequence-level probabilities remain relatively consistent. Thus, GSPO directly addresses and mitigates the volatility associated with expert activation in MoE architectures, rendering workaround strategies such as Routing Replay unnecessary. This leads not only to a more streamlined and reliable training regimen but also enables the MoE model to realize its maximum potential unconstrained by artificial limitations.

\subsection{Benefit of GSPO for RL Infrastructure}

Due to the inconsistencies in precision between training engines (such as Megatron) and inference engines (like SGLang and vLLM), practitioners often rely on the training engine to recalculate the likelihoods of generated responses according to the old policy $\pi_{\theta_\text{old}}$. Nevertheless, since GSPO optimizes using sequence-level likelihoods instead of token-level ones, and the former is inherently more robust to precision mismatches, GSPO allows for the direct use of likelihoods produced by the inference engine in the optimization process. This obviates the need to re-evaluate with the training engine, offering significant advantages in contexts such as partial rollouts, multi-turn reinforcement learning, and frameworks where training and inference systems are separated.

\section{Conclusion}

We introduce Group Sequence Policy Optimization (GSPO), an innovative reinforcement learning algorithm designed for optimizing large language models. Building on the foundations of importance sampling, GSPO utilizes sequence likelihoods to construct importance ratios and incorporates sequence-level strategies for clipping, rewarding, and optimization. Relative to GRPO, GSPO achieves enhanced stability during training, improved sample efficiency, and superior overall performance. It is particularly well-suited for large-scale reinforcement learning training of MoE models and underpins the significant advancements seen in recent Qwen3 models. As a robust and scalable algorithmic framework, GSPO paves the way for ongoing scaling of RL, opening the door to substantial progress in artificial intelligence.

\end{document}